\chapter{Modelo de datos}
\label{cap:datos}

\begin{intuicion}
El grafo es un \emph{multigrafo dirigido temporal} sobre entidades.
``Multi'' = puede haber varias aristas entre los mismos dos nodos
(por ejemplo, el mismo proceso lanzando otro varias veces). ``Dirigido''
= cada arista tiene origen y destino fijos, no son simétricas.
``Temporal'' = cada arista lleva un \texttt{timestamp} y los algoritmos
explotan ese orden. ``Sobre entidades'' = los nodos no son strings sueltos
sino tuplas tipadas \texttt{(EntityType, identifier)}.
\end{intuicion}

\section{\texttt{Entity}, \texttt{Relation}, \texttt{StrId}}

\begin{lstlisting}[language=Rust,caption={Tipos públicos del modelo. Los campos opcionales se omiten en el ejemplo para foco.}]
pub struct Entity {
    pub id: String,                 // ej. "10.0.0.42" o "lsass.exe"
    pub entity_type: EntityType,    // ej. IP, Process, User
    pub metadata: HashMap<...>,     // campos opcionales del log
}

pub struct Relation {
    pub source_id: String,
    pub dest_id: String,
    pub rel_type: RelationType,     // ej. Auth, Connect, Spawn
    pub timestamp: i64,             // segundos epoch
    pub metadata: HashMap<...>,
}
\end{lstlisting}

\texttt{Entity::id} y \texttt{Relation::source\_id}/\texttt{dest\_id}
son \texttt{String}s humanos. Internamente, el motor los pasa por un
\emph{interner} (\texttt{StringInterner}) que devuelve un
\texttt{StrId} (\texttt{u32}) compacto. La \texttt{StringInterner}
guarda cada string una sola vez en un \texttt{Vec<Arc<str>>};
resoluciones futuras devuelven \texttt{Arc<str>} sin reasignación.

\begin{intuicion}
Pensá el \emph{interner} como una tabla de hash bidireccional con
identidad estable. Cada string nuevo entra con un \texttt{StrId}
secuencial; el motor pasa \texttt{StrId}s en lugar de \texttt{String}s
durante todo el \emph{matching} (4~bytes vs.\ ~24~bytes/punter +
\emph{heap} en cada copia). Sólo al \emph{emit} de un resultado se
resuelven a \texttt{Arc<str>}.
\end{intuicion}

\section{\texttt{CompactRelation}: 32~bytes/arista}

El \texttt{Relation} público (humano) pesa $\sim 200$~bytes en RAM.
Para soportar EVTX de 28~GB ($\sim 100$M aristas) sin reventar la
\emph{heap}, el motor usa un \texttt{CompactRelation} de 32~bytes:

\begin{lstlisting}[language=Rust]
#[repr(C)]
pub struct CompactRelation {
    pub timestamp: i64,           // 8 B
    pub metadata_offset: u64,     // 8 B (offset al MetadataStore)
    pub source_sid: StrId,        // 4 B
    pub dest_sid: StrId,          // 4 B
    pub ingested_at_delta: u32,   // 4 B (segundos desde INGEST_EPOCH_BASE)
    pub dataset_tag: u16,         // 2 B
    pub rel_type_tag: u8,         // 1 B
    // 1 B padding implícito
}
\end{lstlisting}

\texttt{metadata\_offset} es un \emph{handle} al
\texttt{MetadataStore} \emph{append-only}; \texttt{0} significa
``sin metadata'' y se chequea sin \emph{lookup}. El \texttt{u8} de
\texttt{rel\_type\_tag} permite el \emph{matching} sin resolver el
\texttt{enum} entero.

\begin{figure}[H]
\centering
\begin{tikzpicture}[
    node distance=0.4cm,
    every node/.style={font=\footnotesize, align=center},
    field/.style={draw=accent, thick, minimum width=1cm, minimum height=0.8cm,
                  fill=accentlight}
]
\node[field, minimum width=2.2cm] (ts) {\texttt{timestamp}\\\scriptsize 8\,B};
\node[field, right=of ts, minimum width=2.4cm] (mo) {\texttt{metadata\_offset}\\\scriptsize 8\,B};
\node[field, right=of mo, minimum width=1.4cm] (src) {\texttt{src\_sid}\\\scriptsize 4\,B};
\node[field, right=of src, minimum width=1.4cm] (dst) {\texttt{dst\_sid}\\\scriptsize 4\,B};
\node[field, below=0.6cm of ts.south west, anchor=north west, minimum width=1.6cm] (ia) {\texttt{ingested\_at}\\\scriptsize 4\,B};
\node[field, right=of ia, minimum width=1.0cm] (dt) {\texttt{ds\_tag}\\\scriptsize 2\,B};
\node[field, right=of dt, minimum width=1.0cm] (rt) {\texttt{rel\_tag}\\\scriptsize 1\,B};
\node[field, right=of rt, minimum width=0.9cm, fill=white, dashed] (pad) {pad\\\scriptsize 1\,B};
\end{tikzpicture}
\caption{\texttt{CompactRelation} en \texttt{\#[repr(C)]}: 32~bytes
exactos, alineados a 8. El \emph{layout} se eligió para que un
\texttt{\&[StreamEdge]} (que comparte \emph{layout} byte-a-byte) se
mapee directamente a una región \emph{mmap} cuando el \emph{streaming
backend} \emph{spills} a disco.}
\label{fig:datos-compact}
\end{figure}

\section{Almacenamiento de aristas: \emph{streaming backend} con auto-spill}

El motor mantiene las aristas en un \emph{streaming backend}
per-vertex \emph{degree-adaptive}: \texttt{Inline} para grado $\leq 8$
(en la struct misma, sin alloc), \texttt{Vec<StreamEdge>} para grado
mayor. Cuando el budget global cruza el umbral
(default $10^6$ aristas en \texttt{GraphHunter::new()},
configurable vía \texttt{with\_memory\_budget\_edges}) el backend
escribe un archivo \emph{round-numbered}, lo \emph{mmap}-ea, y
swap-ea cada \texttt{NeighborList} a la variante \texttt{Mmap \{ handle,
byte\_offset, len \}}. El \emph{sort} por \texttt{timestamp} se hace
\emph{durante el write}, así que el archivo es \emph{matcher-ready}.
Las búsquedas binarias \texttt{partition\_point} sobre la \emph{slice}
funcionan indistintamente sobre la región \emph{Inline}, \emph{Vec}
en RAM, o \emph{Mmap} de disco.

\begin{enelcodigo}
\filepath{core/graph-engine/src/streaming.rs}:
\texttt{InMemoryStreamingBackend} con \texttt{append\_edge},
\texttt{finalize}, \texttt{auto\_spill}, \texttt{neighbors\_arc},
\texttt{for\_each\_edge}. Stage 2c part 2 eliminó el antiguo
\texttt{SpillableEdgeStore} (heavy-logs~\#1/\#2) en favor de este
diseño unificado --- el matcher \emph{hot path} y los $\sim$50
\emph{readers} \emph{admin/aggregation} comparten un solo \emph{store}.
\end{enelcodigo}

\section{\texttt{HuntResult} y la migración a \texttt{Arc<str>}}

Un \emph{hunt} produce \texttt{Vec<HuntResult>}, donde
\texttt{HuntResult = Vec<Arc<str>>}. Cada elemento del \texttt{Vec}
es un vértice del \emph{path} resuelto vía el \emph{interner}.

\begin{intuicion}
El cambio de \texttt{Vec<String>} a \texttt{Vec<Arc<str>>} (heavy-logs)
parece menor pero ahorró $\sim 55\%$ del tiempo de \emph{cap-saturated
benches}. \texttt{String::from(...)} \emph{aloca} bytes en \emph{heap};
\texttt{Arc::clone(...)} sólo bumpea un contador atómico. En un
\emph{hunt} con 10K resultados de 4~vértices cada uno, eso son 40K
\emph{allocs} convertidos en 40K \emph{refcount bumps}.
\end{intuicion}

\section{Índices auxiliares}

Sobre el grafo construido el motor mantiene dos índices que aceleran
el \emph{matcher} pero no son la \emph{source of truth}:

\begin{itemize}
    \item \textbf{NLF} (\emph{Neighborhood Label Frequency}): para
    cada vértice, cuántos vecinos salientes tiene de cada
    \texttt{EntityType}. Permite descartar \emph{starts} cuyo
    histograma no satisface el primer paso de la hipótesis.
    \item \textbf{k-hop reachability}: bitmap de los tipos
    alcanzables en a lo sumo \texttt{MAX\_DEPTH=4} saltos, más
    \texttt{min\_first} y \texttt{max\_first} para gating temporal.
\end{itemize}

Ambos se construyen \emph{lazy} en el primer \emph{hunt} sobre la
sesión (heavy-logs~\#3.5) y se invalidan / actualizan
incrementalmente en cada \texttt{add\_relation} (L3 stage~2a).
Capítulo~\ref{cap:indexes} entra en detalle.
